\documentclass[12pt]{article}
\usepackage{amsmath, amssymb}

\begin{document}

\title{Lecture L15--L18: Joint Random Variables + Joint Distributions}
\date{}
\maketitle

\section{Outline of the Lecture}

The lecture is structured as follows:

\begin{itemize}
\item Why Study Joint Random Variables
\item Discrete Case
\begin{itemize}
\item Joint PMF Table Representation
\item Example: Rolling Two Dice
\end{itemize}
\item Continuous Case
\begin{itemize}
\item Geometric Interpretation
\item Worked Example
\end{itemize}
\item Joint Cumulative Distribution Function
\begin{itemize}
\item Definition of Joint CDF
\item Continuous Case
\item Discrete Case
\item Properties of the Joint CDF
\end{itemize}
\item Comprehensive Example
\end{itemize}

\section{Why Study Joint Random Variables?}

\subsection{Core Statement}

In science and in real life, we are often interested in two (or more) random variables at the same time.

\subsection{Examples}

The lecture provides the following examples:

\begin{itemize}
\item Height and weight of giraffes
\item IQ and birthweight of children
\item Frequency of exercise and rate of heart disease
\item Air pollution level and respiratory illness rate
\item Number of Facebook friends and age of members
\end{itemize}

\subsection{Example \#1: Smart Transportation}

\textbf{Definition of Variables}
\begin{itemize}
\item $X =$ Vehicle Speed
\item $Y =$ Traffic Density
\end{itemize}

\textbf{Question}

Can speed be high when the traffic density is also high?

\textbf{Answer}

Yes / No

\textbf{Statement}

$X$ and $Y$ are correlated random variables

\textbf{Applications (as listed)}

\begin{itemize}
\item Adaptive traffic signals
\item Congestion prediction
\item Autonomous vehicle routing
\end{itemize}

\subsection{Example \#2: Wireless Network Performance}

\textbf{Definition of Variables}
\begin{itemize}
\item $X =$ Signal-to-Noise Ratio (SNR)
\item $Y =$ Data Throughput
\end{itemize}

\textbf{Statement}

Throughput depends on SNR

\textbf{Question}

If SNR is high, will throughput always be high?

\textbf{Answer}

Yes / No

\textbf{Reasoning}

Network congestion (condition relationships)

\textbf{Applications}

\begin{itemize}
\item 5G scheduling
\item Adaptive modulation
\item Network planning
\end{itemize}

\subsection{Example \#3, \#4, \#5}

\begin{itemize}
\item Weather Prediction:
\begin{itemize}
\item $X =$ Rainfall
\item $Y =$ Temperature
\end{itemize}
\item Drone Delivery:
\begin{itemize}
\item $X =$ Wind Speed
\item $Y =$ Delivery Time
\end{itemize}
\item Rolling two dice:
\begin{itemize}
\item $(X,Y)$
\item Conditions such as $X + Y = 7$ and $X > Y$
\item This requires joint PMF
\end{itemize}
\end{itemize}

\subsection{Key Statements}

\begin{itemize}
\item Random variables have a joint distribution
\item This allows:
\begin{itemize}
\item Computing probabilities of events involving both variables
\item Understanding the relationship between the variables
\end{itemize}
\item Special case:
\begin{itemize}
\item Simplest when variables are independent
\end{itemize}
\item When not independent:
\begin{itemize}
\item Use covariance
\item Use correlation
\end{itemize}
\end{itemize}

\section{Discrete Case -- Joint PMF}

\subsection{Setup}

Suppose $X$ and $Y$ are discrete random variables:

\[
X \in \{x_1, x_2, \ldots, x_n\}
\]

\[
Y \in \{y_1, y_2, \ldots, y_m\}
\]

The ordered pair $(X,Y)$ takes values in the product set:

\[
\{(x_1,y_1), (x_1,y_2), \ldots, (x_n,y_m)\}
\]

\subsection{Definition: Joint PMF}

\[
p(x_i, y_j) = P(X = x_i, Y = y_j)
\]

This gives the probability of the joint outcome occurring simultaneously.

\section{Joint PMF Table Representation}

\subsection{Table Organization}

\begin{itemize}
\item Rows correspond to values $x_i$
\item Columns correspond to values $y_j$
\item Entries correspond to $p(x_i, y_j)$
\end{itemize}

\subsection{Key Properties}

\textbf{Property 1}
\[
0 \le p(x_i, y_j) \le 1
\]

\textbf{Property 2}
\[
\sum_{i=1}^{n} \sum_{j=1}^{m} p(x_i, y_j) = 1
\]

\section{Example: Rolling Two Dice}

\subsection{Experiment}

Roll two fair dice.

\subsection{Definition of Random Variables}

\begin{itemize}
\item $X =$ value on the first die
\item $T =$ total (sum) of both dice
\end{itemize}

\subsection{Joint Probability Construction}

Since each outcome $(i,j)$ has probability:

\[
\frac{1}{36}
\]

We construct the joint probability table of $(X,T)$.

\subsection{Observations}

\begin{itemize}
\item Entries are either $0$ or $\frac{1}{36}$
\item $X$ and $T$ are not independent
\end{itemize}

\section{Continuous Case -- Joint PDF}

\subsection{Analogy to Discrete Case}

\begin{itemize}
\item Discrete values $\rightarrow$ Continuous intervals
\item Joint PMF $\rightarrow$ Joint PDF
\item Sums $\rightarrow$ Double integrals
\end{itemize}

\subsection{Domain of $(X,Y)$}

If:
\[
X \in [a,b], \quad Y \in [c,d]
\]

Then:
\[
(X,Y) \in [a,b] \times [c,d]
\]

\subsection{Definition: Joint PDF}

A function $f(x,y)$ is called the joint probability density function of $X$ and $Y$ if:

\[
P((X,Y)\text{ lies in a small rectangle of area } dx,dy) \approx f(x,y)\,dx\,dy
\]

\subsection{Key Properties}

\textbf{Property 1}
\[
f(x,y) \ge 0
\]

\textbf{Property 2}
\[
\int_c^d \int_a^b f(x,y)\, dx\, dy = 1
\]

The joint PDF $f(x,y)$ may take values greater than 1. It is a probability density, not a probability.

\section{Worked Example}

\subsection{Given}

\[
f(x,y) = 4xy, \quad 0 \le x \le 1, \; 0 \le y \le 1
\]

Jointly, $(X,Y)$ takes values in the unit square $[0,1]^2$.

\subsection{Tasks}

\begin{enumerate}
\item Show that $f(x,y)$ is a valid joint PDF
\item Visualize the event:
\[
A = \{X < 0.5, \; Y > 0.5\}
\]
\item Compute $P(A)$
\end{enumerate}

\subsection{Step 1: Check Validity (Normalization)}

\[
\int_0^1 \int_0^1 4xy\, dx\, dy
\]

\[
= \int_0^1 \left(\int_0^1 4xy\, dx\right) dy
\]

\[
= \int_0^1 2y\, dy
\]

\[
= [y^2]_0^1 = 1
\]

\textbf{Conclusion}

\begin{itemize}
\item $f(x,y) \ge 0$
\item Total probability is 1
\item Hence, it is a valid joint PDF
\end{itemize}

\subsection{Step 2: Compute $P(A)$}

\[
P(A) = \int_0^{0.5} \int_{0.5}^{1} 4xy\, dy\, dx
\]

\[
= \int_0^{0.5} \left[2x y^2\right]_{0.5}^{1} dx
\]

\[
= \int_0^{0.5} 2x(1 - 0.25)\, dx
\]

\[
= \int_0^{0.5} \frac{3}{2}x\, dx
\]

\[
= \frac{3}{16}
\]

\section{Joint Cumulative Distribution Function (Joint CDF)}

\subsection{Notation}

\[
X \le x, \; Y \le y
\]

means:

\[
\{X \le x \text{ and } Y \le y\}
\]

\subsection{Definition}

\[
F(x,y) = P(X \le x, \; Y \le y)
\]

\section{Joint CDF -- Continuous Case}

\subsection{Definition via Integration}

\[
F(x,y) = \int_c^y \int_a^x f(u,v)\, du\, dv
\]

\subsection{Recovering Joint PDF}

\[
f(x,y) = \frac{\partial^2}{\partial x \partial y} F(x,y)
\]

\section{Joint CDF -- Discrete Case}

\subsection{Definition}

\[
F(x,y) = \sum_{x_i \le x} \sum_{y_j \le y} p(x_i, y_j)
\]

\subsection{Interpretation}

\begin{itemize}
\item Add all probabilities in the table
\item For rows with $x_i \le x$
\item For columns with $y_j \le y$
\end{itemize}

\section{Properties of the Joint CDF}

\textbf{Property 1: Non-decreasing}

\[
F(x,y) \text{ is non-decreasing}
\]

If $x$ or $y$ increase, then $F(x,y)$ must stay constant or increase.

\textbf{Property 2: Lower-left Corner}

\[
F(x,y) = 0 \text{ at the lower-left of the joint range}
\]

If the lower-left corner is $(-\infty, -\infty)$:

\[
\lim_{(x,y)\to(-\infty,-\infty)} F(x,y) = 0
\]

\textbf{Property 3: Upper-right Corner}

\[
F(x,y) = 1 \text{ at the upper-right of the joint range}
\]

If the upper-right corner is $(\infty, \infty)$:

\[
\lim_{(x,y)\to(\infty,\infty)} F(x,y) = 1
\]

\end{document}